\chapter{CƠ SỞ LÝ THUYẾT}
\label{chap:theory}

Chương này tóm tắt các khái niệm và phương pháp được sử dụng trong đề tài, làm nền tảng để
hiện thực và diễn giải kết quả ở Chương 4.

\section{Tín hiệu giọng nói và các nhóm đặc trưng}

Bộ dữ liệu sử dụng các đại lượng âm học thu được từ bản ghi âm đọc nguyên âm ``aa'' dài
khoảng một giây của từng chủ thể \cite{little2009}. Mỗi bản ghi được mô tả bởi 22 đặc
trưng thuộc năm nhóm chức năng, phản ánh các triệu chứng giọng nói của bệnh Parkinson.

\begin{table}[H]
\centering
\caption{Các nhóm đặc trưng trong bộ dữ liệu Parkinson}
\begin{tabular}{|l|p{7.4cm}|l|}
\hline
\textbf{Nhóm} & \textbf{Đặc trưng đại diện} & \textbf{Số lượng} \\
\hline
Tần số cơ bản & MDVP:Fo(Hz), MDVP:Fhi(Hz), MDVP:Flo(Hz) & 3 \\
\hline
Độ run tần số (Jitter) & MDVP:Jitter(\%), MDVP:Jitter(Abs), MDVP:RAP, MDVP:PPQ, Jitter:DDP & 5 \\
\hline
Độ run biên độ (Shimmer) & MDVP:Shimmer, MDVP:Shimmer(dB), Shimmer:APQ3, Shimmer:APQ5, MDVP:APQ, Shimmer:DDA & 6 \\
\hline
Tỷ lệ nhiễu -- giọng & NHR, HNR & 2 \\
\hline
Chỉ số phi tuyến & RPDE, DFA, spread1, spread2, D2, PPE & 6 \\
\hline
\end{tabular}
\end{table}

Nhóm tần số đoạn hình phổ cơ bản; jitter và shimmer đo dao động ngắn hạn của chu kỳ, thường
tăng ở giọng run; NHR/HNR đo nhiễu tương đối và độ hòa hợp; các chỉ số phi tuyến (RPDE, DFA,
spread1, spread2, D2, PPE) mô tả độ phức tạp và động học không gian trạng thái của tín hiệu
\cite{little2009,han2011}.

\section{Tiền xử lý dữ liệu}

\subsection{Kiểm tra chất lượng}
Các kiểm tra cơ bản bao gồm kích thước, kiểu dữ liệu, giá trị thiếu, bản ghi trùng lặp, giá
trị ngoại lai và độ lệch (skewness). Giá trị ngoại lai được xác định theo quy tắc IQR:

\begin{equation}
\text{IQR} = Q_3 - Q_1, \qquad
\text{ngưỡng} = [Q_1 - 1{,}5\,\text{IQR},\; Q_3 + 1{,}5\,\text{IQR}].
\label{eq:iqr}
\end{equation}

Độ lệch được đo bằng hệ số skewness; với phân phối lệch mạnh ($|\text{skew}| > 1$) cần lưu ý
khi sử dụng các mô hình nhạy cảm với phân phối.

\subsection{Chuẩn hoá}
Vì các đặc trưng có thang đo khác nhau (từ $10^{-6}$ đến hàng trăm Hz), nhóm áp dụng chuẩn
hoá z-score:

\begin{equation}
z_{ij} = \frac{x_{ij} - \mu_j}{\sigma_j},
\label{eq:zscore}
\end{equation}

trong đó $\mu_j$ và $\sigma_j$ lần lượt là kỳ vọng và độ lệch chuẩn của đặc trưng thứ $j$.
Quan trọng, bộ chuẩn hoá chỉ được học trên tập huấn luyện trong các mô hình có giám sát, nhằm
tránh rò rỉ dữ liệu; với PCA và phân cụm, nhóm chuẩn hoá trên toàn bộ tập đặc trưng do không
sử dụng nhãn trong quá trình học.

\section{Phân tích thành phần chính (PCA)}

PCA tìm một hệ trục tuyến tính mới sao cho thành phần đầu tiên giữ lượng phương sai lớn
nhất, các thành phần tiếp theo kề nó trực giao và giữ dần lượt phương sai còn lại
\cite{jolliffe2002}. Gọi $\mathbf{X} \in \mathbb{R}^{n\times p}$ là ma trận đã chuẩn hoá, ta tính
ma trận hiệp phương sai $\mathbf{S} = \frac{1}{n-1}\mathbf{X}^{\mathsf T}\mathbf{X}$, phân tích
thành trị riêng:

\begin{equation}
\mathbf{S}\mathbf{v}_i = \lambda_i \mathbf{v}_i, \quad
\lambda_1 \ge \lambda_2 \ge \dots \ge \lambda_p \ge 0,
\label{eq:eigen}
\end{equation}

và tỷ lệ phương sai giải thích của thành phần thứ $i$ là
$\lambda_i / \sum_{j=1}^{p}\lambda_j$. PCA được dùng ở đây vừa để trực quan hoá dữ liệu trong
không gian giảm chiều, vừa để xác định số thành phần cần giữ đạt ngưỡng phương sai mong muốn.

\section{Các phương pháp phân loại có giám sát}

\subsection{K-Nearest Neighbors (KNN)}
KNN phân loại một điểm bằng cách xem xét $k$ điểm láng giềng gần nhất trong tập huấn luyện
và lấy đa số nhãn. Phương pháp không có bước huấn luyện tham số nhưng nhạy cảm với thang đo
đặc trưng, vì vậy dữ liệu phải được chuẩn hoá.

\subsection{Support Vector Machine (SVM)}
SVM tìm siêu phẳng phân chia tối đu định hướng, tối đu khoảng cách biên (margin) giữa hai lớp;
với ranh giới phi tuyến, ta sử dụng kernel. Trong thực nghiệm nhóm dùng kernel RBF, cho phép
mô hình các ranh giới phi tuyến tổng quát hơn \cite{han2011,james2013}.

\subsection{Random Forest (RF)}
RF là mô hình tập hợp (ensemble) gồm nhiều cây quyết định được huấn luyện trên các mẫu ngẫu
nhiên và tập con đặc trưng ngẫu nhiên, sau đó tổng hợp bằng đa số phiếu. Mô hình thường có
độ chính xác cao nhưng dễ quá khớp (overfit) nếu không giới hạn độ sâu hoặc số lượng mẫu tối
thiểu trên lá \cite{james2013}.

\subsection{Kiểm định chéo và độ đo F-score}
Để ước lượng hiệu năng khách quan, nhóm dùng kiểm định chéo phân tầng (stratified k-fold): tập
dữ liệu được chia $k$ phần giữ nguyên tỷ lệ nhãn, mỗi lần lấy một phần làm kiểm tra và các
phần còn lại làm huấn luyện. Tham số mô hình được chọn bằng tìm lưới (GridSearchCV) trên một
kiểm định chéo trong, tách biệt với kiểm tra ngoài. Với nhãn mất cân bằng, độ đo F-score được
dùng làm độ đo chính:

\begin{equation}
F_1 = \frac{2\,P\,R}{P+R}, \quad
\text{macro-F} = \frac{1}{C}\sum_{c=1}^{C} F_{1,c}, \quad
\text{weighted-F} = \sum_{c=1}^{C} w_c F_{1,c},
\label{eq:f1}
\end{equation}

với $P,R$ lần lượt là precision và recall, $C$ là số lớp, $w_c$ là tỷ trọng mẫu của lớp $c$.

\section{Các phương pháp phân cụm không giám sát}

\subsection{K-Means}
K-Means chia $n$ điểm thành $K$ cụm bằng cách tối thiểu tổng bình phương khoảng cách từ
điểm tới tâm cụm (inertia):

\begin{equation}
J = \sum_{c=1}^{K} \sum_{\mathbf{x} \in C_c} \|\mathbf{x} - \boldsymbol{\mu}_c\|^2.
\label{eq:kmeans}
\end{equation}

Số cụm $K$ được chọn dựa trên chỉ số silhouette trung bình:

\begin{equation}
s(i) = \frac{b(i) - a(i)}{\max\{a(i), b(i)\}},
\label{eq:silhouette}
\end{equation}

trong đó $a(i)$ là khoảng cách trung bình của điểm $i$ tới các điểm cùng cụm, $b(i)$ là
khoảng cách trung bình tới cụm gần nhất; $s(i)$ càng gần 1 tức cụm càng gắn kết và tách biệt.

\subsection{DBSCAN}
DBSCAN nhóm điểm theo mật độ: một điểm là lõi (core) nếu số điểm trong lân cận bán kính
$\varepsilon$ không ít hơn \texttt{min\_samples}; các điểm lõi nối được tạo thành cụm, điểm
không thuộc cụm nào được coi là nhiễu (noise) \cite{ester1996}. DBSCAN không cần khai báo số
cụm và phát hiện được cụm hình dạng tùy ý, nhưng nhạy cảm với $\varepsilon$ và
\texttt{min\_samples}.

\subsection{Chỉ số Rand hiệu chỉnh (ARI)}
ARI đo độ tương đồng giữa phân hoạch cụm và nhãn thật, hiệu chỉnh theo ngẫu nhiên:

\begin{equation}
\text{ARI} = \frac{\text{RI} - \mathbb{E}[\text{RI}]}{\max(\text{RI}) - \mathbb{E}[\text{RI}]},
\label{eq:ari}
\end{equation}

với RI là chỉ số Rand dựa trên số cặp mẫu được phân loại đồng ý. ARI $= 1$ là phù hợp hoàn
hảo, $0$ tương đương với phân ngẫu nhiên, và giá trị âm cho thấy phân hoạch kém hơn ngẫu
nhiên.

\section{Đảm bảo tái lập và chống rò rỉ dữ liệu}

Ba nguyên tắc được áp dụng xuyên suốt đề tài:
\begin{enumerate}
    \item \textbf{Một đầu vào chuẩn duy nhất:} mọi bước phân tích đều đọc từ
          \texttt{parkinsons\_v1.csv} (GATE \#1), tránh lệch phiên bản dữ liệu.
    \item \textbf{Bộ chuẩn hoá nằm trong pipeline:} trong phân loại, \texttt{StandardScaler}
          được fit trong \texttt{Pipeline} nên mỗi fold huấn luyện chỉ học trên phần huấn luyện.
    \item \textbf{Hạt giống cố định:} seed \texttt{42} được đặt ở mọi thủ tục ngẫu nhiên
          (chia dữ liệu, K-Means, Random Forest, GridSearchCV), giúp kết quả tái lập.
\end{enumerate}
